\subsection*{3.6 Inference from Sample Statistics and Margin of Error }

\vspace{0.5em}

\setcounter{questionnum}{0}


\noindent\markforreview

\begin{center}
    

\textbf{Texting Behavior and Cell Phone Use}

\vspace{0.4em}

\begin{tabular}{|l|c|c|c|}
\hline
Texting behavior & Talks on cell phone daily & Does not talk on cell phone daily & Total \\
\hline
Light & 110 & 146 & 256 \\
Medium & 139 & 164 & 303 \\
Heavy & 166 & 74 & 240 \\
\hline
Total & 415 & 384 & 799 \\
\hline
\end{tabular}
\end{center}

\vspace{1em}

In a study of cell phone use, 799 randomly selected US teens were asked how often they talked on a cell phone and about their texting behavior.  
The data are summarized in the table above.  
Based on the data from the study, an estimate of the percent of US teens who are heavy texters is 30\% and the associated margin of error is 3\%.  
Which of the following is a correct statement based on the given margin of error?

\vspace{0.75em}

\choicebox{A}{Approximately 3\% of the teens in the study who are classified as heavy texters are not really heavy texters.}  
\choicebox{B}{It is not possible that the percent of all US teens who are heavy texters is less than 27\%.}  
\choicebox{C}{The percent of all US teens who are heavy texters is 33\%.}  
\choicebox{D}{It is doubtful that the percent of all US teens who are heavy texters is 35\%.}

\vspace{2em}


\newpage
\noindent\markforreview

\begin{center}
    

\textbf{Views on Nuclear Energy Use}

\vspace{0.4em}

\begin{tabular}{|l|c|}
\hline
Response & Frequency \\
\hline
Strongly favor & 56 \\
Somewhat favor & 214 \\
Somewhat oppose & 104 \\
Strongly oppose & 37 \\
\hline
\end{tabular}
\end{center}

\vspace{1em}

A researcher interviewed 411 randomly selected US residents and asked about their views on the use of nuclear energy.  
The table above summarizes the responses of the interviewees.  
If the population of the United States was 300 million when the survey was given, based on the sample data for the 411 US residents,  
what is the best estimate, in \textbf{millions}, of the difference between the number of US residents who somewhat favor or strongly favor  
the use of nuclear energy and the number of those who somewhat oppose or strongly oppose it?  
(Round your answer to the nearest whole number.)

\vspace{0.75em}

\vspace{2em}




\noindent\markforreview


\begin{center}
    

\textbf{Sample A vs. Sample B}

\vspace{0.4em}
\begin{tabular}{|l|c|c|}
\hline
Sample & Percent in favor & Margin of error \\
\hline
A & 52\% & 4.2\% \\
B & 48\% & 1.6\% \\
\hline
\end{tabular}

\end{center}

\vspace{1em}

The results of two random samples of votes for a proposition are shown above.  
The samples were selected from the same population, and the margins of error were calculated using the same method.  
Which of the following is the most appropriate reason that the margin of error for sample A is greater than the margin of error for sample B?

\vspace{0.75em}

\choicebox{A}{Sample A had a smaller number of votes that could not be recorded.}  
\choicebox{B}{Sample A had a higher percent of favorable responses.}  
\choicebox{C}{Sample A had a larger sample size.}  
\choicebox{D}{Sample A had a smaller sample size.}

\vspace{2em}

\newpage
\noindent\markforreview
\vspace{1em}

In State X, Mr. Camp’s eighth-grade class consisting of 26 students was surveyed and 34.6 percent of the students reported that they had at least two siblings.  
The average eighth-grade class size in the state is 26.  
If the students in Mr. Camp’s class are representative of students in the state’s eighth-grade classes and there are 1,800 eighth-grade classes in the state,  
which of the following best estimates the number of eighth-grade students in the state who have fewer than two siblings?

\vspace{0.75em}

\choicebox{A}{16,200}  
\choicebox{B}{23,400}  
\choicebox{C}{30,600}  
\choicebox{D}{46,800}

\vspace{2em}


\noindent\markforreview
\begin{center}
    

\textbf{Poll Results}

\vspace{0.4em}
\begin{tabular}{|l|c|}
\hline
Candidate & Votes \\
\hline
Angel Cruz & 483 \\
Terry Smith & 320 \\
\hline
\end{tabular}
\end{center}

\vspace{1em}

The table shows the results of a poll. A total of 803 voters selected at random were asked which candidate they would vote for in the upcoming election.  
According to the poll, if 6,424 people vote in the election, by how many votes would Angel Cruz be expected to win?

\vspace{0.75em}

\choicebox{A}{163}  
\choicebox{B}{1,304}  
\choicebox{C}{3,864}  
\choicebox{D}{5,621}

\vspace{2em}

\newpage

























\newpage
